% Buck-Boost Converter
% This DC-DC converter could decrease or increase the voltage level from input to output

% Designed by: Amir Ostadrahimi

\documentclass [border=5pt]{standalone}
\usepackage{tikz}
\usepackage[american,cuteinductors,smartlabels]{circuitikz}

%-- the dimensions of the elements can be changed here
\ctikzset{bipoles/thickness=0.7}
\ctikzset{bipoles/length=1.5cm}
\ctikzset{bipoles/resistor/width=0.7}
\ctikzset{bipoles/resistor/height=0.25}
\ctikzset{bipoles/diode/height=0.3}
\ctikzset{bipoles/diode/width=0.3}
\ctikzset{tripoles/thickness=0.7}
\ctikzset{bipoles/battery1/height=0.7}

%settings for fonts and lines
\tikzstyle{every node}=[font=\Large]
\tikzstyle{every path}=[line width=0.9 pt, line cap=round, line join=round]




\begin{document}

 \begin{circuitikz}

%------ Converter
\coordinate  (VBottom) at (0,0); 

\draw (VBottom) to [battery1, l=$V_{in}$, invert] ++(0,3) coordinate (VTop); 

\draw (VTop) to [short] ++(0.8,0) node [nigbt, anchor=C, rotate=90, label={[yshift=-1.5 cm] \normalsize $IGBT$}] (nigbt1){}; 

\draw (nigbt1.E) to [short]++(0.5,0) coordinate (InductorTop) to [short]++(0.5,0) to [D*, invert] ++(2,0) coordinate (DiodeRight); 
 
\draw (InductorTop)  to [L, l_=$L$, v^=$v_L$, i>_=$i_L$] (InductorTop|-VBottom);  

\draw (DiodeRight) to [short]++(0.5,0)  coordinate (CapacitorTop) to [C, l_=$C$, i>^=$i_c$](CapacitorTop|-VBottom); 

\draw  (CapacitorTop) to [short]++(1.5,0)  coordinate (ResistorTop) to [R, l_=$R$, v^=$V_{out}$](ResistorTop|-VBottom) coordinate (ResistorBottom)

(ResistorBottom)-- (VBottom);


%------ Conversion Ratio

\coordinate [label={ [xshift=0, yshift=0] \large $ M(D)=\frac{V_{out}}{V_{in}}=\frac{-D}{1-D};$ \large $0 \leq D <1$ }] (M) at (3,-1);

%%------ Curve

\def\xo{2} % Axes of origin
\def\yo{-2} % Axes of origin
\def\length {5} % length of the axes
\def\N{200} %number of samples

\coordinate [label={ [xshift=-5, yshift=-12]  \normalsize$0$ }] (origin) at (\xo,\yo); %

\draw[->] (\xo -0.2, \yo) --++(\length,0) node[right] {\small $D$};
 \draw[->]  (\xo,\yo +0.2) -- ++(0, -\length) node[below] {\small $M(D)$};
\draw     plot[samples=\N, domain=  0     :  0.8 , xshift=\xo cm, yshift=\yo cm] (5*\x  ,{-\x/(1-\x)});

\draw (\xo , \yo)++(2.5,0) node [above] (half){\normalsize $0.5$};
\draw (\xo , \yo)++(0,-1) node [left] (minus-one){\normalsize $-1$};
\draw [dashed] (half)|-(minus-one);

\end{circuitikz}

\end{document}